\section{Gaps and Limitations}
\label{sec:limitations}

The current artifact is a small, synthetic mechanism-validation study. Its limitations are organized by the failure modes that determine whether the reference architecture can be tested at higher fidelity.

\subsection{Experimental scale}

The pilot uses a shared NumPy multilayer perceptron, a generated scenario grammar, three seeds, and a small sealed split. It is not a trained Transformer, a human-judgment benchmark, a population study, or a deployment evaluation. The task labels are internally generated and may reward feature shortcuts. The next scale-up must use a matched small open-weight Transformer, held-out structural compositions, human-reviewed semantic relations, more seeds, and preregistered stopping and success criteria.

\subsection{Semantic coverage}

The architecture assumes that the relevant patients, relations, norms, exceptions, purposes, and effects can be represented and retrieved. Missing or contradictory structure can produce a well-formed but wrong action. The current coverage test removes generated features; it does not measure coverage against an independently adjudicated world model. A future study needs an explicit coverage audit, missing-context labels, completeness probes, calibrated abstention, and causal tests for whether the missing relation changes the action.

\subsection{Formalization}

The pilot verifier is a transparent rule-based fixture. It is neither an SMT solver nor a proof of semantic completeness. A formal verifier can establish properties of a compiled action abstraction and constraint set; it cannot establish that the abstraction captures every relevant effect, that evidence is true, or that a normative source is legitimate. Future assurance work should compare a fixed action compiler with SMT constraints and state-transition specifications such as TLA+ while reporting unknown and out-of-model cases.

\subsection{Normative grounding}

The L0/L1 division, the candidate value vocabulary, authority metadata, and conflict states are design choices. Signatures, provenance, and versioning provide traceability but do not determine who has legitimate authority over a person or jurisdiction. Pluralistic evaluation can preserve disagreement and profile adherence; it cannot settle moral or political legitimacy. Normative late binding also creates a governance surface for source selection, scope, effective dates, exceptions, appeal, and rollback.

\subsection{Neural representation}

The compiler $C_\phi$ may lose relation semantics, entangle values with surface features, or create an apparently stable shortcut. Auxiliary objectives can interfere with general capability, calibration, or language understanding. Structural attention and routing may improve a benchmark by exploiting graph regularities rather than the intended concepts. Relation permutation, random-label controls, concept-direction probes, representation drift tests, and matched-compute baselines are required before claims about semantic grounding.

\subsection{Security}

Explicit semantic inputs add attack surfaces: source and authority spoofing, ontology poisoning, purpose manipulation, semantic occlusion, routing evasion, prompt injection, world-state poisoning, compiler compromise, and verifier bypass. The current attacks are synthetic fixtures and cover only the paths encoded by the generator. The threat matrix below records candidate controls and residual risks.

\section{Alignment and security threat model}
\label{sec:threats}

Explicit semantic structure may improve observability and constrain some behaviors, but it also creates attack surfaces. An attacker might poison the ontology, falsify world-state evidence, manipulate routing, exploit conflicts between norms, or bypass the verifier. The threat model is therefore part of the architecture rather than an afterthought.

\begin{table*}[t]
\centering
\scriptsize
\caption{Threat classes, affected layers, and proposed controls. Residual risk remains an empirical question.}
\label{tab:threats}
\begin{tabularx}{\textwidth}{@{}p{0.16\textwidth}p{0.13\textwidth}X X p{0.17\textwidth}@{}}
\toprule
Threat & Attacked layer & Failure mode & Candidate mitigation & Residual risk \\
\midrule
Prompt injection / jailbreak & Runtime context & Adversarial text overrides relevant norms or purpose & Separate authority channels, adversarial training, semantic conflict detection & Novel encodings and ambiguous authority \\
Goal misgeneralization & Neural computation & Capability transfers while intended objective changes & OOD goals, counterfactual tests, structured salience supervision & Unseen causal structure \\
Reward hacking & Training objective & Proxy optimization scores well while violating intended purpose & Multiple metrics, gold audits, value-structure probes & Proxy gaps remain \\
Ontology poisoning & L0/L1/L2 & Malicious relation or definition changes decisions & Signed versions, provenance, review, regression and rollback & Governance capture \\
World-state poisoning & L3 & False or stale evidence drives a permitted action & Source trust, temporal validity, confidence, cross-checking & Correlated evidence failure \\
Normative drift / axiology corruption & L0/L1 & Continued learning changes protected semantics & Frozen or gated artifacts, drift tests, approval boundary & Representation drift without schema change \\
Routing manipulation & Neural computation & Context causes relevant experts to be skipped & Router audits, permutation tests, fallback coverage & Adversarial distribution shift \\
Verifier bypass & Assurance boundary & Unstructured or alternate tool path escapes checking & Enforced action schema and policy enforcement point & Implementation defects \\
Specification gaming & Constraint layer & Formally compliant action is materially harmful & Independent review, richer effects model, red-team tests & Incomplete formalization \\
Conflicting ontologies & L1/L2 & Authoritative sources disagree and model hides conflict & Explicit conflict state, precedence metadata, calibrated abstention & No universally accepted resolution \\
Governance compromise & All layers & Canonical value source is captured or malicious & Separation of duties, public change logs, rollback and review & Institutional legitimacy \\
\bottomrule
\end{tabularx}
\end{table*}

Prompt injection and jailbreaks test whether context can overwhelm learned or structural constraints \citep{zou2023universal,wei2023jailbroken}. Goal misgeneralization and reward hacking test whether a system pursues a proxy or an unintended objective while retaining competence \citep{langosco2022goalmisgeneralization,shah2022goalmisgeneralization,gao2023overoptimization,denison2024subterfuge}. Ontology and world-state poisoning add a distinct concern: a model may reason correctly from a compromised symbolic substrate. The appropriate response is not to assume that explicit structure is trustworthy, but to make provenance, authority, uncertainty, and rollback first-class inputs.

Specification gaming is especially important for the formal assurance layer. A verifier can only check the representation and constraints it receives. If the neural system omits a moral patient, chooses a misleading abstraction, or finds an action that satisfies a narrow predicate while producing a harmful effect, a solver can return a valid result for the wrong problem. Independent semantic review and post-deployment monitoring are therefore necessary.

The security objective is not “no attacks.” It is measurable reduction in attack success, faster detection, bounded blast radius, preserved provenance, and safe failure when the architecture is uncertain or contradictory.


\subsection{Scalability and operations}

Runtime context assembly, provenance retention, conflict handling, graph retrieval, structural conditioning, routing, action compilation, verification, and audit storage add latency and operational complexity. Approximate FLOPs are recorded for the current CPU-only proxy, but the artifact does not yet measure Transformer memory, latency, throughput, energy, or production rollback behavior. A larger study must report compute-normalized gains, failure blast radius, version compatibility, and cost of assurance.
